\documentclass[11pt,letterpaper]{article}
\usepackage[T1]{fontenc}
\usepackage[utf8]{inputenc}
\usepackage{lmodern,microtype}
\usepackage[margin=0.94in,headheight=15pt]{geometry}
\usepackage{amsmath,amssymb,amsthm,mathtools}
\usepackage{booktabs,longtable,array,enumitem,fancyhdr,xcolor,listings,needspace}
\usepackage[hyphens]{url}
\usepackage{hyperref,bookmark}
\definecolor{ink}{RGB}{22,45,65}
\definecolor{pale}{RGB}{246,248,250}
\hypersetup{colorlinks=true,linkcolor=ink,urlcolor=blue!45!black,citecolor=ink,
 pdftitle={Counter Congruence Suffices for iPRA},
 pdfauthor={Research prepared for Vladimir Reshetnikov}}
\setlength{\emergencystretch}{3em}
\setlength{\parskip}{0.25em}
\setlist{itemsep=0.25em,topsep=0.4em,leftmargin=1.6em}
\newtheorem{theorem}{Theorem}[section]
\newtheorem{lemma}[theorem]{Lemma}
\newtheorem{proposition}[theorem]{Proposition}
\newtheorem{corollary}[theorem]{Corollary}
\theoremstyle{definition}\newtheorem{definition}[theorem]{Definition}
\theoremstyle{remark}\newtheorem{remark}[theorem]{Remark}
\newcommand{\ipra}{\mathsf{iPRA}}
\newcommand{\Acal}{\mathcal K_{\mathrm{ctr}}}
\newcommand{\D}{\ensuremath{\mathsf D}}
\newcommand{\Q}{\ensuremath{\mathsf Q}}
\newcommand{\av}{\mathrel{\triangleright}}
\newcommand{\FV}{\operatorname{FV}}
\newcommand{\llbracket}{[\![}
\newcommand{\rrbracket}{]\!]}
\newcommand{\file}[1]{{\small\nolinkurl{#1}}}
\lstset{basicstyle=\ttfamily\footnotesize,columns=fullflexible,keepspaces=true,
 breaklines=true,showstringspaces=false,frame=single,rulecolor=\color{black!15},
 backgroundcolor=\color{pale},xleftmargin=0.3em,xrightmargin=0.3em}
\pagestyle{fancy}\fancyhf{}
\fancyhead[L]{\small Counter congruence suffices for iPRA}
\fancyhead[R]{\small Results and implementation limits}
\fancyfoot[C]{\thepage}
\title{\color{ink}\textbf{Counter Congruence Suffices for iPRA}\\[0.5em]
 \large A structural reduction of equality rules,\\
 with obstructions to two seedless recursion bases}
\author{Research prepared for Vladimir Reshetnikov}
\date{November 10, 2026}
\begin{document}
\maketitle
\begin{abstract}
With the existing successor-seeded, counter-aware recursor, congruence in
its counter alone, together with replacement at the left root of an
equation, derives congruence through every fixed numerical expression.
The construction uses numerical cases, not equation induction, and keeps
both numerical computation schemas unchanged. An inherited lifting then
recovers arbitrary formula transport. The resulting complete native
calculus is contextually equivalent to the existing iPRA presentation;
its whole-domain, actual-equality bi-interpretation and equiconsistency
are given explicitly.

The requested significant checker-complexity breakthrough is \emph{not}
established. The inspected corpus's coordinate minima are 281 normalized
Python tokens, 315 AST nodes, and 685 normalized bytes, attained by
different complete programs. Incidental, behavior-preserving Python
changes improve these coordinates separately to 280, 314, and 672.
They implement the existing theory and do not give one program below
all three old minima. The new restricted calculus has no size record.
Its positive contribution is fixed-counter sufficiency under the
unmodified recursor equations. The numerical construction specializes
an existing compiler, and the older initializer strategy is acknowledged.

Two additional theorems show why simply deleting the independent
initializer is inadequate for direct term representations: zero-seeded
counter recursion preserves monotonicity, while a seemingly stronger
successor-seeded variant admits a unary bound depending on only one
input coordinate. A third audit identifies a witness-order obstruction
to a naive reversal of the paired binder's quantifiers.
\end{abstract}
\begin{center}
\fcolorbox{ink}{pale}{\begin{minipage}{0.91\linewidth}
\textbf{Status.} No significant new axiomatization below the complete
checker frontier, or absolute simplestness, is established. The report
supplies a complete equivalent calculus, a restriction theorem, small
implementation records for the inherited theory, executable certificates,
and scoped obstructions. No Lean/Rocq formalization or external human
peer review is claimed.
\end{minipage}}
\end{center}
\clearpage
\setcounter{tocdepth}{2}
\begingroup\small\setlength{\parskip}{0pt}\tableofcontents\endgroup
\clearpage

\section{The benchmark and the precise contribution}
\label{sec:benchmark}

The inspected ProveIt snapshot is commit
\begin{center}\file{2fdcf5bf988bead821c416649f713f98b7ed8776}.\end{center}
There are nine direct ZIP files in \file{docs/drafts}. Each was retrieved
and checked against both its byte length and Git blob identity. A
recursive inventory, including the historical archives embedded inside
the current packages, contains 14 distinct ZIP contents, 342 Python
source occurrences, and 164 distinct Python source hashes. The inventory
and the manual distinction between complete checkers and auxiliary
scripts are included in \file{evidence/}.

\begin{center}\small
\begin{tabular}{lrrr}
\toprule
Entire complete program & Tokens & AST nodes & Bytes\\
\midrule
Generated assumptions, principal &281&318&685\\
Generated assumptions, AST profile &283&315&710\\
Guarded elimination, token profile &281&325&739\\
Guarded elimination, successor byte profile &290&333&689\\
Atomic transport, unrestricted control &297&335&698\\
\midrule
Inspected-corpus coordinate envelope &\textbf{281}&\textbf{315}&\textbf{685}\\
\bottomrule
\end{tabular}
\end{center}

The measurements use the unchanged inherited normalizer under CPython
3.12.14: remove comments and docstrings, apply \texttt{ast.unparse}, and
count significant tokens, AST nodes, and UTF-8 bytes of the whole
module. Input handling, scope checks, every inference branch, and the
public verdict driver count. The host Python runtime and standard JSON
library are common infrastructure. Reference-directory checkers count
when they are complete programs. Tiny build and test wrappers do not.

The older 287/323/738 envelope printed in several drafts is obsolete
for the snapshot examined here. Beating that older envelope would not
meet the user's current target. Section~\ref{sec:incidental} records
small implementation-only improvements found during the final audit;
they change the coordinate envelope to 280/314/672. These coordinates
still belong to different programs. Moreover all these attained sizes are
\emph{upper bounds} on achievable checker size. They are not a theorem
giving an absolute lower bound, which would require a fixed description
language and cost convention.

\subsection{What changes mathematically}

General equality transport is normally presented with an arbitrary
formula or numerical context as rule data. Here it is recovered from two
much more specific operations: replacing the complete left endpoint of
an equation, and changing only a recursor's counter while keeping its
initializer and update fixed. The replacement of a rule by two more
specific rules is a restriction of its primitive interface; it does not
by itself imply shorter checker source.

The useful observation is that every fixed numerical context can be
represented, provably, by two nested recursors in whose varying arguments
the hole occurs only as a counter. This makes the restricted congruence
rule sufficient. The proof is given before any appeal to the full iPRA
interpretation, so it does not assume the equality principles being
reconstructed.

\subsection{Attribution and novelty limits}

The successor-seeded recursor, paired binder, case/progress distinction,
logical macros, ordinary arithmetic bridge, and lifting of atomic to
compound transport are inherited\cite{atomic,successor,fixedposition}.
The context construction below is a specialization of the existing
primitive-recursion compiler. We do not claim a new recursor, a new
compiler for all primitive recursive functions, or a new compound
transport theorem.

The older fixed-position report already uses root replacement and
replacement in the \emph{initializer} of a scalar iterator. Its direct
context representation uses an additional first-step computation schema
I1. Here the counter-aware R basis is unchanged: R0 and RS suffice, and
the only congruence position needed is the \emph{counter}. No matching
restricted-rule theorem was found in the current and nested textual
sources or in the separately retrieved full historical report. This is
a novelty statement relative to the inspected corpus, not a claim of
priority over the literature.

\section{Complete primitive syntax and rules}
\label{sec:theory}

\subsection{Numerical terms and computation}

Terms consist of variables, zero, and the counter-aware recursion form
\[
                       R(n,a;u,v.h).
\]
Only the update body $h$ binds the two distinct variables $u,v$; its
remaining parameters stay fixed through the recursion. Write
\[
 S(a):=R(0,a;u,v.0),\qquad 1:=S0.
\]
For readability, the implementation represents $S$ as an explicit
constructor. Its defining equality with the displayed R term is already
an R0 instance; it adds no independent arithmetic axiom.

The only computation schemas, available under any retained context, are
\begin{align}
 R(0,a;u,v.h)&=S(a),\tag{R0}\label{eq:R0}\\
 R(Sn,a;u,v.h)&=h[n/u,R(n,a;u,v.h)/v].\tag{RS}\label{eq:RS}
\end{align}
Substitutions are capture-free. The recursive state retains the original
update $h$. Each body is a fixed finite term description. Equivalently,
introduce a first-order symbol $r_h$ at each finite syntactic stage.
There is no object-level variable ranging over functions or programs,
and no assertion of a uniformly total primitive-recursive interpreter.

\subsection{Object formulas and internal clauses}

Object formulas are
\[
                         A::=(s=t)\mid\D_{x,y}A,
\]
where $x,y$ are distinct bound numerical variables. A context is a finite
reusable list of object formulas and internal clauses $\Q_yB(y)$, with
$B$ an object formula. The ordinary readings are
\[
 \D_{x,y}B\;\mapsto\;
   \exists x\bigl(B(x,0)\to\forall y\,B(x,y)\bigr),
 \qquad
 \Q_yB\;\mapsto\;B(0)\to\forall y\,B(y).
\]
An internal clause is available only by context membership; it is not
an exported object goal. An object formula is available either by
membership or by a proof. Write this availability as $\Gamma\av M$.
The primitive hypothesis rule is
\[
                \frac{}{\Gamma\vdash A}\quad(A\in\Gamma).
                \tag{Hyp}
\]
Here and below $A$ as a proof conclusion is an object formula. The
designated contradiction is $\bot_*=(0=1)$.

\subsection{Restricted equality rules}

\[
 \frac{\Gamma\av s=t\qquad\Gamma\vdash s=r}
      {\Gamma\vdash t=r}
 \quad\mathsf{RL},
 \tag{Root replacement}\label{rule:RL}
\]
\[
 \frac{\Gamma\av s=t}
      {\Gamma\vdash R(s,a;u,v.h)=R(t,a;u,v.h)}
 \quad\mathsf{CC}.
 \tag{Counter congruence}\label{rule:CC}
\]
CC has exactly the same initializer, update, and passive parameters on
both sides, modulo renaming of bound variables. It gives no primitive
permission to alter an initializer, a successor argument, an update,
or a compound formula. Reflexivity, symmetry, transitivity, and every
other congruence are derived.

\subsection{All remaining logical rules}

\[
 \frac{\Gamma\av\Q_yB(y)\qquad\Gamma\vdash B(0)}
      {\Gamma\vdash B(t)}\quad\mathsf{QE},
 \tag{Closure use}\label{rule:QE}
\]
\[
 \frac{\Gamma\av\D_{x,y}B(x,y)\qquad
       \Gamma,\Q_yB(a,y)\vdash F}
      {\Gamma\vdash F}\quad\mathsf{DE},
 \tag{Paired elimination}\label{rule:DE}
\]
\[
 \frac{\Gamma,B(t,0)\vdash B(t,Sz)}
      {\Gamma\vdash\D_{x,y}B(x,y)}\quad\mathsf{DI}_{\rm case},
 \tag{Case introduction}\label{rule:DIcase}
\]
\[
 \frac{\Gamma,E(t,z)\vdash E(t,Sz)}
      {\Gamma\vdash\D_{x,y}E(x,y)}\quad\mathsf{DI}_{\rm progress}.
 \tag{Equation progress}\label{rule:DIprogress}
\]
In DE, $a$ is fresh for $\Gamma,F$ and for all fixed parameters of the
matrix. In both introduction rules, $z$ is fresh for $\Gamma,t$ and
the fixed matrix parameters. Progress is permitted only when $E$ is
a \emph{literal numerical equation}. A formula with an equation under
a paired binder is compound and cannot use progress. Local hypotheses
and clauses are discharged. Bound variables are renamed before all
substitutions.

These rules, Hyp, R0, and RS are the entire primitive inventory. There
is no primitive general equality transport, coverage axiom, separation
axiom, general formula induction, excluded middle, or unlisted ambient
first-order proof checker. Weakening, exchange, contraction, admissible
renaming, numerical substitution, and object cut follow by external
induction on finite derivations. The usual freshening of local
parameters makes substitution and cut preserve all side conditions.
Denote the calculus by $\Acal$.

% Inclusion fragment. Requires amsmath, amssymb, amsthm and theorem,
% lemma, proposition, corollary environments. No private macros are used.
% This fragment matches Eq/D object formulas, internal Q clauses, and
% ordinary hypothesis use for object formulas only.

\section{Counter congruence recovers unrestricted equality transport}
\label{ctr:section}

The reduction in this section concerns the mathematical inference
interface. A smaller primitive interface and a smaller executable checker
are separate claims; only the accompanying measurements can establish
the latter. We use the numerical terms, the Eq/$\mathsf D$ object
formulas, the internal $\mathsf Q$ clauses, and the primitive rules
specified above. In particular,
\[
 \begin{aligned}
 R(0,a;u,v.h)&=Sa,\\
 R(Sn,a;u,v.h)&=h[n/u,R(n,a;u,v.h)/v],
 \end{aligned}
 \tag{R0, RS}
\]
where $Sa:=R(0,a;u,v.0)$.
Only the update $h$ binds $u,v$, and all substitutions are capture-free.
No variable ranges over numerical program descriptions. A template in
what follows is one fixed finite expression.

Recall the two restricted equality rules:
\[
 \frac{\Gamma\mathrel{\triangleright}s=t\qquad
       \Gamma\vdash s=r}
      {\Gamma\vdash t=r}
 \quad\mathsf{RL},
 \qquad
 \frac{\Gamma\mathrel{\triangleright}s=t}
      {\Gamma\vdash R(s,a;u,v.h)=R(t,a;u,v.h)}
 \quad\mathsf{CC}.
 \tag{Root and counter}
\]
In $\mathsf{CC}$ the seed, update, and passive parameters are the same
on both sides. The availability relation $\Gamma\mathrel{\triangleright}M$
means that $M$ is a retained assumption or an already proved object
formula. Internal $\mathsf Q$ clauses are available by membership.
Full hypothesis use is present for Eq/$\mathsf D$ object assumptions;
there is no rule with an internal $\mathsf Q$ clause as its object goal.

The proof uses $\mathsf{Hyp}$, $\mathsf{QE}$, $\mathsf{DE}$ and
$\mathsf{DI}$-case, in addition to R0, RS, $\mathsf{RL}$ and
$\mathsf{CC}$. It does not use $\mathsf{DI}$-progress. In particular,
no equation-induction instance is hidden in the reduction.

\subsection{Structural substitution is not equality transport}

Weakening, exchange, contraction, renaming, and capture-free substitution
of a numerical term for a free parameter are admissible by induction on
finite derivations. Rename eigenparameters first. R0 and RS are stable
under substitution; substitution preserves the matching seed and update
in $\mathsf{CC}$ and the matching root in $\mathsf{RL}$. It also
preserves a formula's root constructor, so a permitted progress instance
remains literal if progress occurs elsewhere in the derivation.
The ordinary eigenparameter checks handle the logical rules.
Object-assumption cut is obtained by replacing each use of the assumption
by the supplied derivation, again freshening eigenparameters.
These are operations on fixed finite proof syntax. They do not appeal
to an object-theory rule replacing equal terms inside a formula.

\subsection{Equality paths using root replacement alone}

\begin{lemma}[Equality algebra]
\label{ctr:equality-algebra}
Reflexivity, symmetry, and transitivity are derivable using RS and
$\mathsf{RL}$, without $\mathsf{CC}$ or general equality transport.
\end{lemma}
\begin{proof}
For a numerical term $s$ put
\[
             c_s:=R(Ss,0;u,v.u).
\]
RS gives $c_s=s$. Use one copy as the available major and another as
$\mathsf{RL}$'s minor; the conclusion is $s=s$.
For symmetry, an available $s=t$ and the just-derived $s=s$ give
$t=s$ by $\mathsf{RL}$.
For transitivity, suppose $s=t$ and $t=r$ are available or proved.
Hypothesis use supplies an object proof when necessary. Symmetry gives
$t=s$; apply $\mathsf{RL}$ with this major and minor $t=r$ to obtain
$s=r$. Consequently every finite path of proved numerical equations
can be composed, and any edge may be reversed. This path operation
changes no proper numerical subterm.
\end{proof}

\subsection{Numerical cases before congruence}

\begin{lemma}[Numerical cases]
\label{ctr:cases}
For any fixed Eq/$\mathsf D$ object template $A(y)$, the following rule
is admissible without either equality rule or progress introduction:
\[
 \frac{\Gamma\vdash A(0)\qquad\Gamma\vdash A(Sz)}
      {\Gamma\vdash A(t)}.
 \tag{Cases}
\]
Here $z$ is fresh for the retained context, the fixed parameters of
$A$, and the target term $t$; it may always be renamed to meet these
conditions.
\end{lemma}
\begin{proof}
Choose a bound coordinate $x$ absent from $A$. By weakening, the
successor premise remains available under the extra assumption
$A(0)$. Apply $\mathsf{DI}$-case with witness zero to obtain
$\Gamma\vdash\mathsf D_{x,y}A(y)$. Eliminate this paired formula.
Because $x$ is unused, its new internal clause is exactly
$\mathsf Q_yA(y)$; the fresh existential parameter does not occur in
that clause. Weaken the given base proof into the extended context and
apply $\mathsf{QE}$ at $t$. This proves $A(t)$ there, and
$\mathsf{DE}$ discharges the clause. There was no rewrite by an
equality and no induction hypothesis.
\end{proof}

\subsection{Representing a numerical context through the counter}

Define
\[
                  W(n):=R(n,0;u,v.u).
\]
The computations $W(0)=S0$ and $W(Sn)=n$ are instances of R0 and RS.
In particular the latter identity is available before any general
congruence principle has been established.

\begin{lemma}[Counter representation of a fixed context]
\label{ctr:context-representation}
For any fixed numerical template $F(x)$, choose all new bound variables
fresh for its passive parameters and put
\[
 \begin{aligned}
  H_F(n)&:=R(n,F(0);u,v.S(F(Su))),\\
  J_F(n)&:=W(H_F(n)).
 \end{aligned}
 \tag{Counter representation}
\]
Then the restricted calculus proves
\[
                         J_F(n)=F(n)
 \tag{Representation}\label{ctr:representation-equation}
\]
for arbitrary numerical $n$ and arbitrary retained context. The proof
uses numerical cases and no progress introduction.
\end{lemma}
\begin{proof}
The update of $H_F$ does not contain its state variable $v$.
Consequently R0 and RS give the following identities literally, with
no equality substitution into an enclosing expression:
\[
 H_F(0)=S(F(0)),\qquad H_F(Sz)=S(F(Sz)).
 \tag{Two computations}
\]
For either of these instances, write its argument as $q$, where $q=0$
or $q=Sz$. Apply $\mathsf{CC}$ to the displayed equality using seed
zero and the update that returns its old counter. This produces
\[
                 W(H_F(q))=W(S(F(q))).
\]
The right-hand expression has the computation equation
$W(S(F(q)))=F(q)$. Compose the two equation paths by
Lemma~\ref{ctr:equality-algebra}; the result is
$J_F(q)=F(q)$. We have now supplied both premises of
Lemma~\ref{ctr:cases} for the literal equation
$J_F(y)=F(y)$. Cases proves
\eqref{ctr:representation-equation} at an arbitrary $n$.

Every proper-subterm change in this argument is the single, explicitly
stated $\mathsf{CC}$ instance through $W$. The remaining equation
compositions use root replacement. In particular one must not silently
replace $H_F(q)$ by $S(F(q))$ under $W$ before applying
$\mathsf{CC}$; that is precisely the congruence being justified.
\end{proof}

The displayed $J_F$ is a specialization of the successor-recursion
compiler already present in the inspected drafts: take its initializer
to be $F(0)$ and its step to be $h(u,z)=F(Su)$, which ignores the state.
The additional observation here is that this specialization is proved
correct using only root replacement and congruence in the counter.
The general recursion compiler itself is inherited.

\begin{theorem}[All numerical congruence from counter congruence]
\label{ctr:all-congruence}
For every fixed numerical template $F(x)$, the restricted calculus admits
\[
 \frac{\Gamma\mathrel{\triangleright}s=t}
      {\Gamma\vdash F(s)=F(t)}.
 \tag{Numerical congruence}
\]
\end{theorem}
\begin{proof}
Use the major $s=t$ in $\mathsf{CC}$ with seed $F(0)$ and update
$S(F(Su))$. Its conclusion is $H_F(s)=H_F(t)$. A second
$\mathsf{CC}$ instance, now for $W$, gives
$J_F(s)=J_F(t)$. Lemma~\ref{ctr:context-representation} supplies
$J_F(s)=F(s)$ and $J_F(t)=F(t)$. Symmetry and path composition from
Lemma~\ref{ctr:equality-algebra} give
\[
                F(s)=J_F(s)=J_F(t)=F(t).
\]
These are equations joined at their roots. Neither the proof nor the
statement assumes an interpretation limited to standard numerals.
\end{proof}

\begin{corollary}[Arbitrary atomic transport]
\label{ctr:atomic-transport}
For all fixed numerical templates $U(x),V(x)$,
\[
 \frac{\Gamma\mathrel{\triangleright}s=t\qquad
       \Gamma\vdash U(s)=V(s)}
      {\Gamma\vdash U(t)=V(t)}
 \tag{Atomic transport}
\]
is admissible.
\end{corollary}
\begin{proof}
Theorem~\ref{ctr:all-congruence} gives $U(s)=U(t)$ and $V(s)=V(t)$.
Reverse the first equality and compose only at equation roots:
\[
                U(t)=U(s)=V(s)=V(t).
\]
\end{proof}

\subsection{A constant bound on inference-node overhead}

Count one node for each primitive inference, including a hypothesis
lookup, and count a proof tree with repeated subproofs repeated. In the
explicit expansion supplied with the report, reflexivity has three
nodes. Symmetry adds four nodes to its supplied proof, and transitivity
adds five nodes to the sum of its two supplied proofs.

Each of the zero and successor proofs of $J_F(q)=F(q)$ has eight
nodes: a computation, one CC lifting, a second computation, and the
five-node transitivity overhead. The cases construction consists of
DE, DI-case, QE, one clause hypothesis lookup, and these two eight-node
proofs. Thus $J_F(t)=F(t)$ has twenty nodes.

If the supplied proof of $s=t$ has $P$ nodes, the two CC liftings have
$P+2$ nodes, the reversed left representation has $20+4$ nodes,
and the right representation has twenty. Two transitivity constructions
therefore give
\[
                  (P+2)+(20+4)+20+5+5=P+56.
\]
The additional inference-node count is independent of the syntactic
depth of the fixed numerical template $F$. Freshening binders changes
names, not the number of inferences. This is not a constant byte-size
or running-time bound: the proof's term and formula labels contain
copies of $F$ and of the supplied endpoints. It is also not a constant
bound for the subsequent compound-formula lifting. The delivered
successor, initializer, and update-body congruence certificates each
have 57 nodes when their supplied equality is a one-node computation.

\subsection{Compound formulas and internal clauses}

The following lifting pattern is inherited from the earlier
fixed-position and atomic-transport arguments. It is included to show
that the numerical result recovers the full logical interface, and to
make its independence from compound induction explicit.

\begin{theorem}[Full object-formula transport]
\label{ctr:full-transport}\label{thm:atomic-transport}
For every fixed Eq/$\mathsf D$ object template $A(w)$, the rule
\[
 \frac{\Gamma\mathrel{\triangleright}s=t\qquad
       \Gamma\vdash A(s)}
      {\Gamma\vdash A(t)}
 \tag{Object transport}
\]
is admissible. Its expansion introduces no progress instances.
\end{theorem}
\begin{proof}
Induct externally on the number of paired binders in $A$. The induction
statement is uniform in all retained contexts, source terms, and passive
parameters. For a literal equation, use
Corollary~\ref{ctr:atomic-transport}.

Suppose $A(w)=\mathsf D_{x,y}B(w,x,y)$. Rename the two bound
variables away from $s,t$ and the retained context. Eliminate the given
$A(s)$ at a fresh parameter $a$. The extended context retains
\[
                     \mathsf Q_y B(s,a,y).
\]
To obtain the target $\mathsf D_{x,y}B(t,x,y)$, use
$\mathsf{DI}$-case with witness $a$ and a fresh successor parameter
$z$. Its local object assumption is $B(t,a,0)$; hypothesis use gives
an object proof of it. Symmetry gives $t=s$. The external induction
hypothesis applied backwards to the smaller template $B(w,a,0)$
therefore gives $B(s,a,0)$. Specialize the retained clause at $Sz$ to
obtain $B(s,a,Sz)$. Apply the induction hypothesis forwards along
$s=t$, now to $B(w,a,Sz)$, obtaining $B(t,a,Sz)$. This is exactly
the required case-introduction premise. Close $\mathsf{DI}$-case and
then $\mathsf{DE}$.

The new parameter $z$ is absent from the retained clause, $a$, and all
fixed matrix parameters; the eliminated witness $a$ is absent from the
outer conclusion. Both recursive transport uses concern the strictly
smaller object matrix $B$. This is recursion on finite external syntax,
not an induction rule of the arithmetic. The only introductions added
in the binder case are case introductions.
\end{proof}

\begin{proposition}[Transporting the use of an internal clause]
\label{ctr:clause-transport}
An available clause $\mathsf Q_yB(s,y)$, together with an available
$s=t$, can be used with exactly the operational force of
$\mathsf Q_yB(t,y)$: from any object proof of $B(t,0)$ one can
obtain $B(t,r)$ for any numerical $r$.
\end{proposition}
\begin{proof}
Transport the supplied base backwards along $t=s$, specialize the old
clause at $r$, and transport the result forwards along $s=t$, using
Theorem~\ref{ctr:full-transport} for the two object matrices.
This statement does not add a rule proving an internal clause as an
object formula. If an object representative is required, use
$\mathsf D_{x,y}B(t,y)$ with unused $x$; case introduction and the
just-described clause use construct that representative.
\end{proof}

\begin{corollary}[Exact identity-language equivalence]
\label{ctr:exact-equivalence}
Let $\mathcal T_{\mathrm{ctr}}$ have R0, RS, object hypothesis use,
$\mathsf{RL}$, $\mathsf{CC}$, $\mathsf{QE}$,
$\mathsf{DE}$, and the case and literal-equation progress introductions.
Let $\mathcal T_{\mathrm{tr}}$ instead use unrestricted Eq/$\mathsf D$
object transport as its equality rule, with all other numerical and
logical syntax and rules the same. For every finite context of object
assumptions and internal clauses, and every object formula $A$,
\[
 \Gamma\vdash_{\mathcal T_{\mathrm{ctr}}}A
 \quad\Longleftrightarrow\quad
 \Gamma\vdash_{\mathcal T_{\mathrm{tr}}}A.
\]
The translations are effective on finite derivations and preserve their
end sequents literally.
\end{corollary}
\begin{proof}
From right to left, translate premise derivations recursively and replace
each unrestricted transport instance by
Theorem~\ref{ctr:full-transport}. All other rules are copied.
From left to right, $\mathsf{RL}$ is the unrestricted-transport instance
with motive $w=r$. Reflexivity is derived from a computation equation
by that same root instance. For $\mathsf{CC}$, begin with reflexivity
$R(s,a;h)=R(s,a;h)$ and transport its left occurrence of $s$ to $t$;
this gives $R(t,a;h)=R(s,a;h)$. Derived symmetry gives the required
orientation. The other rules are again copied. All side conditions are
preserved by freshening local parameters.
\end{proof}

Because this equivalence leaves the native language and all primitive
meanings unchanged, the whole-domain bi-interpretations and actual-equality
composite isomorphisms of $\mathcal T_{\mathrm{tr}}$ transfer to
$\mathcal T_{\mathrm{ctr}}$. A proof of that arithmetic and logical
bridge is supplied separately; exact sequent equivalence is not a
substitute for stating the bridge. The present argument supplies no
absolute lower bound on checker size and no claim of novelty outside
the specifically inspected sources.

% Inclusion fragment. Requires amsmath, amssymb, amsthm and theorem,
% proposition/corollary environments. No private macros are used.

\subsection{Counter congruence alone does not recover root replacement}
\label{rootind:section}

The two restricted equality rules are not interchangeable. The following
counterinterpretation shows that simply deleting root replacement loses
even the theorem $0=0$. This is a statement about the explicitly specified
native inference system. It does not rule out every different rule,
encoding, numerical basis, or interpretation.

\begin{theorem}[Root replacement is not derivable from the remaining rules]
\label{rootind:theorem}
Let $\mathcal T^{-\mathsf{RL}}$ have the same terms, Eq/$\mathsf D$
object formulas, internal $\mathsf Q$ clauses, and rules R0, RS,
$\mathsf{Hyp}$, $\mathsf{CC}$, $\mathsf{QE}$, $\mathsf{DE}$,
$\mathsf{DI}$-case, and $\mathsf{DI}$-progress as the counter-congruence
calculus, but omit $\mathsf{RL}$. Do not add an unlisted reflexivity,
equality-substitution, or other ordinary equality rule. Then
\[
                    \varnothing\nvdash_{\mathcal T^{-\mathsf{RL}}}0=0.
\]
In particular, $\mathsf{RL}$ is not admissible in this reduced system.
\end{theorem}
\begin{proof}
For this syntactic nonderivability argument, interpret the native atom
written $s=t$ by a binary predicate $E$. We are not assuming that this
predicate is actual equality: the missing equality principles are
precisely what the proposed reduction would have to derive.

Take numerical domain $\{0,1\}$. Interpret zero as $0$ and every
recursion expression by
\[
       \llbracket R(n,a;u,v.h)\rrbracket_\rho=1
\]
under every assignment $\rho$, independently of its counter, seed,
update, and passive parameters. This is compatible with capture-free
substitution and alpha-renaming. The successor abbreviation
$Sa=R(0,a;u,v.0)$ is therefore the constant function $1$, and the
numeral $1=S0$ denotes the domain element $1$.
Interpret the equality-shaped atom by
\[
       \llbracket s=t\rrbracket_\rho\text{ is true}
       \quad\Longleftrightarrow\quad
       \llbracket s\rrbracket_\rho=1.
 \tag{Predicate interpretation}
\]
Its right-hand term is immaterial. The equality on the right of this
metatheoretic definition is the usual equality of the two domain elements.
Interpret the binders and internal clauses by their ordinary meanings:
\[
 \begin{aligned}
  (\mathsf Q_y B)^\circ
    &:=B^\circ(0)\to\forall y\,B^\circ(y),\\
  (\mathsf D_{x,y}B)^\circ
    &:=\exists x\bigl(B^\circ(x,0)\to\forall y\,B^\circ(x,y)\bigr).
 \end{aligned}
\]
These are ordinary two-valued semantics, equivalently a one-world
intuitionistic interpretation.

We check every retained primitive rule under arbitrary assignments and
arbitrary true retained contexts.

\paragraph{Computation and counter congruence.}
Each R0 or RS axiom has an $R$ expression as the left side of its atom.
That expression has value $1$, so the atom is true, irrespective of its
right side. A $\mathsf{CC}$ conclusion also has an $R$ expression on
its left and is therefore true. These rules are sound even though the
binary predicate is not reflexive at zero.

\paragraph{Hypotheses, availability, and elimination.}
An object hypothesis is true whenever its retained context is true.
An available major is either such a retained assumption or an already
proved object formula, so availability preserves truth.
For $\mathsf{QE}$, a true clause
$B(0)\to\forall yB(y)$ and a true base $B(0)$ give its requested
instance $B(t)$. For $\mathsf{DE}$, the true existential major supplies
a domain element satisfying the retained closure clause. The fresh
witness parameter can be assigned that element; its absence from the
outer conclusion and retained context gives the ordinary soundness
argument for existential elimination.

\paragraph{Case introduction for arbitrary object matrices.}
Fix an assignment making $\Gamma$ true, and let $b$ be the value of
the submitted witness $t$. Freshness makes $b$ and $\Gamma$
independent of the new successor parameter $z$. A sound premise
\[
        \Gamma,B(t,0)\vdash B(t,Sz)
\]
ensures that $B(b,0)$ implies $B(b,1)$, since $Sz$ always denotes
$1$. The domain has just $0$ and $1$, so this implication ensures
$B(b,0)\to\forall yB(b,y)$. Choosing $x=b$ proves the intended
$\mathsf D$ conclusion. This verifies $\mathsf{DI}$-case without
restricting its matrix to an equation.

\paragraph{Progress introduction for literal equations.}
Fix the passive parameters and the witness. By inspection of finite
numerical term syntax, a term $U(y)$ has one of the following three
value functions of the remaining coordinate $y$:
\[
                y\longmapsto0,\qquad
                y\longmapsto1,\qquad
                y\longmapsto y.
\]
Indeed a zero term is constant $0$; a fixed passive variable is constant
$0$ or $1$; the distinguished variable has value $y$; and any term
with an $R$ root, including a successor abbreviation, is constant $1$.
For a literal equation $E(y)=(U(y)=V(y))$, its truth depends only on
$U(y)$ by the predicate interpretation. Its truth set is consequently
one of
\[
                       \varnothing,\qquad\{1\},\qquad\{0,1\}.
\]
In every case, truth of $E(0)$ implies truth of $E(y)$ for both domain
elements. Thus the required closure
$E(0)\to\forall yE(y)$ is valid for every fixed witness, even before
using the progress premise. This proves soundness of
$\mathsf{DI}$-progress. The restriction to a literal equation is
essential to this verification; it is not a claim about arbitrary
compound matrices.

All primitive rules of $\mathcal T^{-\mathsf{RL}}$ are now sound.
Induction on a finite derivation therefore preserves truth in this
interpretation. But the atom $0=0$ is false because its left term has
value $0$. It has no derivation from the empty context.

Finally, in the full calculus RS gives
\[
                  R(S0,0;u,v.u)=0.
\]
Two copies of this computation, used as the major and minor of
$\mathsf{RL}$, prove $0=0$. If $\mathsf{RL}$ were admissible in
the reduced system, this three-node argument could be translated into
it, contradicting the counterinterpretation.
\end{proof}

The conclusion is deliberately limited. Deleting $\mathsf{RL}$ destroys
the established identity-language equivalence with the unrestricted
transport calculus. The counterinterpretation is not asserted to be a
model of an ordinary theory in which equality already carries its usual
logical axioms, and it is not an impossibility theorem for all
bi-interpretations with a different equality or domain. Importing those
usual equality axioms silently would change the proof-checker problem.


\section{How the equivalence claim is completed}
\label{sec:bridge-guide}

The new restriction theorem proves exact consequence equivalence with
the inherited calculus, using the identity map on numerical terms,
object formulas, and contexts. This gives more than a standard-model
simulation: every finite proof can be translated in both directions.
To make the stronger bi-interpretability claim reviewable, the next
three sections give the arithmetic and logical bridge in full.

Those sections adapt the atomic-transport report at the pinned snapshot,
with attribution\cite{atomic}. Every invocation of general transport
is justified by the preceding theorem; every object assumption is
available by Hyp. The earlier report's elimination-only implementation
is not silently imported. In particular, the two interpretation
composites are checked for arbitrary elements and fixed descriptions,
with equality given by actual numerical equality.

% Adapted with attribution from the atomic-transport draft at the pinned repository snapshot.
\section{The short numerical construction}
\label{sec:numerical}

The following direct compiler is inherited from the successor-recursion
draft\cite{recursor}; it is included to make the present construction
self-contained. Atomic equality reasoning already gives transitivity
and numerical congruence. Define
\[
 W(n):=R(n,0;u,v.u),\qquad
 C(n,a,b):=W(R(n,a;u,v.Sb)),\qquad
 P(n):=C(n,0,W(n)).
\]
Computation and derived equality reasoning prove
\begin{align*}
 W0&=1,&W(Sn)&=n,\\
 C(0,a,b)&=a,&C(Sn,a,b)&=b,\\
 P0&=0,&P(Sn)&=n.
\end{align*}
The counter input $u$ of R is essential: it is not an alias for $v$.
Applying W proves that successor is injective.

Suppose fixed terms $g(\vec p)$ and $h(n,z,\vec p)$ are already
represented. Put
\[
 J_{g,h}(n,\vec p)=
 W\bigl(R(n,g(\vec p);u,v.S h(u,W(v),\vec p))\bigr).
 \tag{Compiler}\label{eq:compiler}
\]
Then
\[
 J_{g,h}(0,\vec p)=g(\vec p),\qquad
 J_{g,h}(Sn,\vec p)=h(n,J_{g,h}(n,\vec p),\vec p).
\]
At zero, R returns $Sg$ and W removes the successor. At a successor,
R returns $S h(n,J_{g,h}(n,\vec p),\vec p)$ and W again removes
the displayed successor. Both identities hold at arbitrary source
elements; there is no standard-numeral restriction or positivity premise.
Zero, successor, and projections are terms; composition is substitution.
External recursion on any fixed PR description therefore compiles it.
No scalar predecessor search, pairing construction, or equation induction
is needed for this compiler. Scalar iteration is its special case with
$h$ independent of its counter.

\section{Deriving the full intuitionistic logical interface}
\label{sec:logic}

This section adapts the inherited Choice-DI interface\cite{elimination,recursor}.
Every retained object formula is available by the Hyp rule, and
every compound equality transport means Theorem~\ref{thm:atomic-transport}.
Thus none of the following reasoning adds a primitive checker rule.

\subsection{Cases, equation induction, and congruence of the binder}

With an unused coordinate $x$, progress DI introduces $\D_{x,y}E(y)$
from $\Gamma,E(z)\vdash E(Sz)$ for a literal equation E. DE opens
its clause $\Q_yE(y)$; a supplied base proof $E(0)$ and QE at t give
$E(t)$. Hence contextual equation induction is admissible:
\[
 \frac{\Gamma\vdash E(0)\qquad\Gamma,E(z)\vdash E(Sz)}
      {\Gamma\vdash E(t)}. \tag{EI}\label{rule:EI}
\]
Case DI in the same construction, ignoring its base assumption in a
given successor proof, yields numerical cases for any object template:
\[
 \frac{\Gamma\vdash A(0)\qquad\Gamma\vdash A(Sz)}
      {\Gamma\vdash A(t)}. \tag{NC}\label{rule:NC}
\]
The usual freshness conditions are retained. Witness 1 may be used
for the unused coordinate. These derivations never need a proof whose
conclusion is an internal Q clause.

Uniform pointwise interderivability of matrices is preserved under D.
Eliminate the source D at a fresh witness, introduce the target D at
that witness by case DI, convert its base backwards, use the opened
source clause at $Sz$, and convert the result forwards. This is the
same constructive pattern as the transport theorem, now with arbitrary
uniform conversion derivations. Reverse the conversions for the converse.

\subsection{Numerical gates and formula selectors}

For a numerical control z define, externally on fixed formula syntax,
\[
 G_z(s=t):=(C(z,0,s)=C(z,0,t)),\qquad
 G_z(\D_{x,y}B):=\D_{x,y}G_zB.
\]
Freshen binders away from the control. Computation, reflexivity, and
binder congruence prove $\vdash G_0A$ and $G_{Sz}A\dashv\vdash A$.
The padding operation $\operatorname{Pad}(A)=\D_{u,y}G_yA$, with
unused u,y, is equivalent to A: case DI proves the forward direction;
DE, the true zero gate, and QE at 1 prove the reverse direction.

Define the selector on two equation leaves by
\[
 M_z(s=t,r=w):=(C(z,s,r)=C(z,t,w)).
\]
At aligned fresh binders extend it by
\begin{align*}
 M_z(\D_{x,y}A,\D_{x,y}B)&=\D_{x,y}M_z(A,B),\\
 M_z(E,\D_{x,y}B)&=\D_{x,y}M_z(G_yE,B),\\
 M_z(\D_{x,y}A,E)&=\D_{x,y}M_z(A,G_yE),
\end{align*}
where E is an equation. The maximum input binder depth decreases in
each recursive call. Padding and binder congruence give
\[
             M_0(A,B)\dashv\vdash A,\qquad
             M_{Sz}(A,B)\dashv\vdash B.
\]
These are finite syntactic operations, not object-language function symbols.

\subsection{Connectives and quantifiers}

Use the scalar iteration I compiled by \eqref{eq:compiler}, put
$Z(n)=C(n,1,0)$ and $p(n)=I(n,0;v.Z(v))$. The only parity facts
needed are the finite computations $p(1)=1$ and $p(2)=0$.
With fresh bound variables and unused u define
\begin{align*}
 A\Rightarrow_*B&=\D_{u,z}M_z(A,B),\\
 A\wedge_*B&=\D_{u,z}G_zM_{p(z)}(A,B),\\
 \exists_*x A(x)&=\D_{x,y}G_yA(x),\\
 A\vee_*B&=\exists_*x M_x(A,B),\\
 \forall_*x A(x)&=A(0)\wedge_*\D_{u,z}G_zA(z),\\
 \bot_*&=(0=1).
\end{align*}

Implication introduction converts its retained zero selector into A,
applies the supplied derivation of B, and converts B to the successor
selector; case DI closes. Elimination is DE followed by QE at 1.
Existential introduction uses its supplied witness and the successor
gate. Elimination opens D, supplies the theorem given by its zero gate,
and uses QE at 1. These are the usual eigenparameter conditions.

Disjunction introductions use control witnesses 0 and 1. For elimination,
turn the two branch derivations into $A\Rightarrow_*F$ and
$B\Rightarrow_*F$. NC gives $M_z(A,B)\Rightarrow_*F$ at any control;
eliminate the existential selector and apply the implication.
For conjunction introduction NC gives $M_v(A,B)$ at every v from the
two supplied proofs. Specialize to $p(Sz)$, add the successor gate,
and use case DI. Projections are QE at 1 and 2 after DE, using the two
displayed parity values.

NC now gives the coverage formula
\[
             (x=0)\vee_*\exists_*y\,(x=Sy).
\]
At zero use the left injection of reflexivity; at a successor use the
right injection with its evident witness. Coverage permits a case split
even if x occurs in retained assumptions: it adds an equality naming
the case, rather than illicitly substituting into the whole context.
Universal introduction obtains the zero and successor cases of a
fresh-parameter proof and forms the two conjuncts. Universal elimination
uses coverage, the first conjunct at zero, and the gated closure at a
successor. This avoids an eigenvariable assumption about the retained context.

Finally, under $0=1$, computation and atomic transport prove
\[
                   s=C(0,s,t)=C(1,s,t)=t
\]
for arbitrary s,t. Every equation follows; external recursion on object
syntax, using case DI at each D, proves every object formula. This
derives explosion for the designated absurdity. We have thus derived
the full intuitionistic natural-deduction rules, with actual equality,
without postulating excluded middle for quantified formulas.

\section{Full arithmetic and both interpretation composites}
\label{sec:interpretations}

\subsection{An explicit ordinary realization}

Let $T_R$ be ordinary intuitionistic first-order logic with equality
in the staged R signature, the equations R0 and RS, contextual
literal-equation induction EI, and
\[
 \forall x\,(x=0\vee\exists y\,x=Sy),\qquad (0=1)\to\bot.
 \tag{Boundary}
\]
For definiteness, ordinary logic here has hypothesis, implication
introduction and modus ponens, conjunction introduction and its two
projections, the two disjunction introductions and case elimination,
universal introduction and specialization, existential introduction and
fresh-witness elimination, and explosion. Equality has reflexivity and
substitution. The usual eigenvariable restrictions are imposed. These
ordinary rules specify the comparison theory, not additional primitives
of $\Acal$.

Realize its object and clause syntax by
\[
 (s=t)^\circ=(s=t),\quad
 (\D_{x,y}B)^\circ=\exists x(B^\circ(x,0)\to\forall yB^\circ(x,y)),
\]
\[
                 (\Q_yB)^\circ=B^\circ(0)\to\forall yB^\circ(y).
\]
The interpretation domain is $x=x$ and equality is actual equality.
Case DI is sound by coverage. Progress DI is sound by equation
induction under its temporarily assumed base. All other native rules
have their direct ordinary meanings. Conversely Section~\ref{sec:logic}
derives the ordinary rules, coverage, and EI. Separation translates to
the identity implication $(0=1)\Rightarrow_*(0=1)$.

The explicit falsity bridge matters. The positive numerical equations
alone admit a singleton structure, so declaring them an ordinary
first-order theory while silently omitting separation would be a mistake.

\subsection{Actual equality reflection}

For clarity, the remaining arithmetic bridge is given internally,
not merely by agreement on standard numerals. Every fixed PR description
is already compiled by Section~\ref{sec:numerical}. In particular use
\[
 x+0=x,\quad x+Sy=S(x+y),\qquad
 d(x,0)=x,\quad d(x,Sy)=P(d(x,y)).
\]
Let $\operatorname{sg}(x)=C(x,0,1)$. Coverage and the computation
equations show that sg takes only 0,1, that
$\operatorname{sg}(x)=0\leftrightarrow x=0$, and $C(x,a,a)=a$.
Applying sg and separation gives $Sx\ne0$. Coverage and W give
$W(x)=0\leftrightarrow x=1$.

Literal-equation induction proves, in order,
\[
 0+y=y,\qquad (Sx)+y=S(x+y),\qquad
 d(Sx,Sy)=d(x,y),\qquad d(Sa,a)=1.
\]
For the first two induct on y with x fixed. For the third induct on y:
the base is $P(Sx)=x$, and its step applies P to the induction equation.
For the fourth induct on a and use the third identity in the step.

One more equation induction, on y with x fixed, gives
\[
        C(\operatorname{sg}(d(x,y)),x,y+d(x,y))=x. \tag{Residual}
\]
The base is $C(\operatorname{sg}(x),x,x)=x$. In the step set
$r=d(x,y)$ and $r'=Pr=d(x,Sy)$. Coverage splits $r'$ into zero
or successor. At zero the selector already returns x. In the successor
case, $r'\ne0$ forces r to be a successor, and $P(Sq)=q$ gives
$r=Sr'$. The induction equation says $y+r=x$. The new second branch
is $(Sy)+r'=y+Sr'=y+r=x$. This completes the induction without
using a decidable actual-equality predicate in its premises.

Consequently $d(Sa,b)=1\leftrightarrow a=b$: forwards, Residual yields
$b+1=Sa$, hence $Sb=Sa$, and W cancels the successors; backwards
use $d(Sa,a)=1$. Thus
\[
                 e(a,b)=\operatorname{sg}(W(d(Sa,b)))
\]
is Boolean and satisfies $e(a,b)=0\leftrightarrow a=b$ for arbitrary
elements of every source model.

For a fixed quantifier-free formula A, define its zero-true
characteristic by
\begin{align*}
 b_{s=t}&=e(s,t),&b_\bot&=1,\\
 b_{A\wedge B}&=C(b_A,b_B,1),&
 b_{A\vee B}&=C(b_A,0,b_B),&
 b_{A\to B}&=C(b_A,b_B,0).
\end{align*}
An external induction on A, using the proved Boolean ranges, gives
$b_A\in\{0,1\}$ and $b_A=0\leftrightarrow A$ intuitionistically.
The implication case uses that $b_A=1$ proves $\neg A$; it does not
assert excluded middle for arbitrary quantified formulas. Applying EI
to $b_A(n)=0$ recovers every contextual quantifier-free induction instance.

\subsection{The two numerical interpretations}

Take iPRA in its full first-order language of fixed PR descriptions,
with defining equations, constructor separation, and quantifier-free
induction. Its atomic-induction formulation is equivalent by the
characteristic argument above; this is the usual first-order iPRA
convention\cite{salehi}. Interpret its descriptions in $T_R$ by
\eqref{eq:compiler}, leaving domain, equality, and ordinary logic unchanged.
All defining equations, successor axioms, and induction instances have
just been verified internally.

Conversely, each fixed R body is interpreted in iPRA by ordinary
primitive recursion with base $Sa$ and step $h(u,v,\vec p)$. The
abbreviation $S(a)=R(0,a;u,v.0)$ then has ordinary successor value.
R0, RS, the boundary sentences, and EI are iPRA theorems.

For the stronger bi-interpretability claim, first anchor the successor
abbreviation. Its iPRA interpretation is the zero instance of the
recursor with zero update, so its defining base clause gives ordinary
$Sa$. Compiling that interpreted successor back gives the original
$S(a)$ by the compiler's base computation. This includes the expansion
$r_0(0,a)$ used to name successor, and does not presuppose uniqueness
of all R symbols. Thus the same successor occurs in both copies of each
recursion before the following staged uniqueness argument begins.

Now fix independent initializer and parameter variables. Two solutions
of the same fixed R recursion
are internally equal by induction on their value equation: bases agree,
and their steps agree by equality congruence and the induction equation.
Induct externally on the finite staged description. Its double
translation satisfies the same equations once earlier subdescriptions
have been identified, so uniqueness identifies it with the original.
The identical argument works for ordinary PR descriptions. Only after
these identities have been proved for independent parameters do we
substitute counter-dependent actual arguments. This also establishes
coherence when substitution changes the name of a staged body.

Both numerical interpretation domains are the entire original domain;
both internally definable isomorphisms have graph $x=y$. Nothing is
restricted to the standard cut or quotiented by a numerical observer.

\subsection{The logical round trips}

The ordinary meanings of gates and selectors give the advertised
meanings of all six macros in Section~\ref{sec:logic}, with coverage
in the universal case and separation for falsity. For the only primitive
object binder the reverse composite is explicitly
\[
 \D_{x,y}B\dashv\vdash
 \exists_*x\bigl(B(x,0)\Rightarrow_*\forall_*yB(x,y)\bigr).
\]
Forwards, open D at fresh a, use its available clause under the assumed
base to prove each $B(a,y)$, form the conditional universal, and introduce
the existential at a. Backwards, eliminate the existential at fresh a,
introduce the original D by case DI, and use its retained base, the
conditional, and universal specialization to obtain the successor case.
Neither direction requires Q to be an object goal. External induction
on formulas establishes both logical composites.

\begin{theorem}[Whole-domain bi-interpretation]
The ordinary realization of $\Acal$ is bi-interpretable with full
first-order iPRA. The two numerical isomorphisms are internally given
by $x=y$, and both logical formula composites are provably equivalent
to their originals. Effective finite proof translations preserve the
designated contradiction, hence the two theories are equiconsistent.
\end{theorem}
This is an external metatheorem. No claim is made that either theory
proves its own consistency or the uniform correctness of an arbitrary
primitive-recursive evaluator.



% Inclusion fragment. Requires amsmath, amssymb, amsthm and theorem,
% lemma, proposition, corollary environments. No private macros are used.

\section{Why two seedless recursion shortcuts do not give a term compiler}
\label{seedless:section}

The two results below concern the standard numerical functions denoted
by fixed finite terms. They do not rule out every possible relational
interpretation, quotient interpretation, or change of domain. In
particular, failure to define a function by a term is weaker than failure
to define its graph and prove its totality. The results nevertheless
exclude the straightforward term-translation strategy proposed for
these two reduced numerical grammars.

\subsection{Zero-seeded counter recursion preserves monotonicity}

Consider numerical terms generated by variables, zero, successor, and
\[
                  L(n;u,v.h),
\]
where only the update $h$ binds $u,v$. Give each fixed update its standard
natural-number interpretation by
\[
 L(0;u,v.h)=0,\qquad
 L(Sn;u,v.h)=h[n/u,L(n;u,v.h)/v].
 \tag{Zero-seeded recursion}
\]
The update may contain passive parameters and smaller recursion terms;
there is no quantified function input.

\begin{proposition}[Monotonicity obstruction]
\label{seedless:monotonicity}
Every standard function denoted by a term of this grammar is
coordinatewise nondecreasing. Therefore its terms cannot denote all
primitive-recursive functions; in particular they cannot denote the
zero indicator $z(0)=1$, $z(Sn)=0$.
\end{proposition}
\begin{proof}
Induct on finite term syntax. Variables, zero, and successor have the
required property. For the recursion constructor, the induction hypothesis
makes the standard update function $h(u,v,\vec p)$ coordinatewise
nondecreasing. With $\vec p$ fixed, write
\[
 r_0=0,\qquad r_{k+1}=h(k,r_k,\vec p).
\]
We have $r_0\le r_1$ because all values are natural numbers. If
$r_k\le r_{k+1}$, monotonicity of $h$ in its first two coordinates gives
\[
 r_{k+1}=h(k,r_k,\vec p)
 \le h(k+1,r_{k+1},\vec p)=r_{k+2}.
\]
Thus the value is nondecreasing in its counter. Separately, induction on
$k$ shows $r_k(\vec p)\le r_k(\vec p\,')$ whenever
$\vec p\le\vec p\,'$ coordinatewise. Combining the two comparisons
proves joint monotonicity in the counter and passive parameters.
Substituting the actual counter term and parameter terms preserves this
property by the structural induction hypothesis. The zero indicator
violates monotonicity at $0<1$.
\end{proof}

\subsection{A more promising seedless operator still has only one unbounded input}

Now consider the grammar with variables, zero, and
$B(n;u,v.h)$, defined by
\[
 \begin{aligned}
 B(0;u,v.h)&=S(h[0/u,0/v]),\\
 B(Sn;u,v.h)&=h[n/u,B(n;u,v.h)/v],\\
 Sa&:=B(0;u,v.a).
 \end{aligned}
 \tag{Successor at the base}\label{seedless:B}
\]
The last line is an abbreviation and agrees with ordinary successor in
the intended standard interpretation. The base value is tied to the
same update that governs the successor steps; there is no independent
initializer argument.

This operator escapes the previous monotonicity obstruction.
For example, the following standard identities hold:
\[
 \begin{array}{lll}
 W(n):=B(n;u,v.u),&W(0)=1,&W(Sn)=n,\\
 Z(n):=B(n;u,v.0),&Z(0)=1,&Z(Sn)=0,\\
 H(n,a):=B(n;u,v.a),&H(0,a)=Sa,&H(Sn,a)=a.
 \end{array}
\]
In particular $P(n):=W(H(n,n))$ is the ordinary truncated predecessor:
$P(0)=0$ and $P(Sn)=n$. Even a data gate is available. Define
\[
        G(n,a):=P\bigl(P(B(a;u,v.H(n,v)))\bigr).
\]
At $n=0$ the internal recursion starts at $2$ and applies successor,
so its value is $a+2$ and $G(0,a)=a$. At $n=Sq$ it starts at $1$
and applies the identity, so $G(Sq,a)=0$. These standard computations
are also formalizable using equation induction in an ordinary
presentation of the displayed equations. They do not, however, supply
a general selector or initialized primitive recursion.

\begin{theorem}[One-coordinate unary bound]
\label{seedless:one-coordinate}
Every fixed term has a bound depending on at most one input.
Precisely, for a term $t(x_1,\ldots,x_k)$ of the displayed $B$ grammar,
one of the following holds in the standard natural numbers:
\begin{enumerate}
 \item $t$ has a constant bound independent of all inputs;
 \item for some coordinate $i$ and some unary primitive-recursive
       function $F$, fixed by the term,
       \[
       t(x_1,\ldots,x_k)\le F(x_i)
       \quad\text{for every }(x_1,\ldots,x_k)\in\mathbb N^k.
       \]
\end{enumerate}
The coordinate and its bounding function may differ between terms.
\end{theorem}
\begin{proof}
Induct on the finite syntax. The variable and zero cases are immediate;
the optional successor abbreviation also preserves the claim. At a
recursion term $B(n;u,v.h)$, the induction hypothesis applies to $h$
as a term in the two distinct bound variables $u,v$ and its passive
parameters. Let $r_j$ denote the standard recursive values at a fixed
assignment of passive parameters.

If $h(u,v,\vec p)\le F(p_i)$ for a passive coordinate, then
\[
 r_0=1+h(0,0,\vec p)\le1+F(p_i),\qquad
 r_{j+1}=h(j,r_j,\vec p)\le F(p_i).
\]
Hence $1+F(p_i)$ is a bound independent of the counter and every other
parameter. A constant bound on $h$ is handled in exactly the same way.

If the controlling coordinate of $h$ is its counter, so
$h(u,v,\vec p)\le F(u)$, define the nondecreasing unary function
\[
                       M(q):=1+\max_{j\le q}F(j).
\]
The base and every successor value at counter $q$ are bounded by
$M(q)$. By the induction hypothesis the actual counter term $n$ is
bounded by a constant or by a unary function $N(x_i)$. Monotonicity of
$M$ gives respectively the constant bound $M(c)$ or the one-coordinate
bound $M(N(x_i))$.

Finally, suppose the controlling coordinate of $h$ is its state:
$h(u,v,\vec p)\le F(v)$. Put
\[
 b:=1+F(0),\qquad
 K(m):=m+1+\max_{j\le m}F(j).
\]
The function $K$ is nondecreasing, dominates $F$, and satisfies
$K(m)\ge m$. Ordinary numerical induction in the metatheory gives
\[
                          r_q\le K^q(b).
\]
Indeed the base is bounded by $b$, and the step follows from
$r_{q+1}\le F(r_q)\le K(r_q)\le K(K^q(b))$.
Since $K\ge\operatorname{id}$, the bound $q\mapsto K^q(b)$ is
nondecreasing. Compose it with the constant or one-coordinate bound
on the actual counter term $n$, as in the preceding case.

Bounded maxima, fixed-function iteration, and composition of fixed
primitive-recursive descriptions are primitive recursive. Thus each
bound used here is a unary primitive-recursive function. These bounding
functions belong to the ordinary metatheory; the argument does not
assert that they are terms of the proposed reduced grammar. Every
possible controlling coordinate of $h$ has been handled, completing
the structural induction.
\end{proof}

\begin{corollary}[Addition and selection]
\label{seedless:no-selector}
The displayed $B$ grammar has no term for ordinary addition
$(a,b)\mapsto a+b$ and no term for the selector
\[
 C(c,a,b)=
 \begin{cases}
 a,&c=0,\\
 b,&c>0.
 \end{cases}
\]
Consequently it cannot provide a term interpretation of the full
primitive-recursive signature by the proposed direct compiler.
\end{corollary}
\begin{proof}
Addition cannot have a bound depending on just $a$: hold $a$ fixed and
let $b$ increase. The symmetric argument excludes a bound depending
only on $b$; a constant bound is also impossible.

For the selector, a bound depending only on $c$ fails at $c=0$ as $a$
increases. A bound depending only on $a$ fails when $a$ is fixed,
$c=1$, and $b$ increases. A bound depending only on $b$ fails when
$b$ is fixed, $c=0$, and $a$ increases. These exhaust its input
coordinates, contradicting Theorem~\ref{seedless:one-coordinate}.
\end{proof}

The obstruction is stronger than an unsuccessful search for a selector.
It applies to every finite term of this grammar, including arbitrary
nesting of its recursion constructor and the displayed predecessor,
zero indicator, and gate as abbreviations. The gate is consistent with
the theorem because it is bounded by its one data argument $a$.
The independent initializer in the established $R$ basis permits the
initial state and the iteration length to carry two independent
unbounded inputs; the above argument explains why deleting that
initializer is a substantive loss for this particular term language.


\section{Witness order in the attempted reversed binder}
\label{sec:reversed}

One tempting simplification replaces the inherited
$\exists x\,\forall y$ order by
\[
 \mathsf E_{x,y}B
       :=\forall y\,\exists x\bigl(B(x,0)\to B(x,y)\bigr).
\]
The hope is to combine the two eliminations and remove the reusable
internal Q clause. Dummy coordinates and numerical gates still suggest
formulas for implication, existence, and universality. That observation
alone is insufficient for bi-interpretability.

Consider the naive case rule that chooses a fixed numerical witness
term $t(z)$ as part of the introduction instruction and then checks
$B(t(z),0)\vdash B(t(z),Sz)$. Its progress variant uses a witness
independent of $z$ and a literal equation. The source existential
witness obtained at a fresh universal parameter may need to be chosen
\emph{inside} that parameter's proof scope, after an existential
elimination. Choosing it before the subproof does not permit this
interleaving.

There is a semantic obstruction to the intended reverse formula
composite. On standard $\mathbb N$, interpret the raw E binder by
\[
 \text{there exists a unary primitive-recursive }f\text{ such that }
 \forall y\,[B(f(y),0)\to B(f(y),y)].
 \tag{PR-selector reading}\label{eq:prselector}
\]
Free parameters are held fixed. The naive case rule is sound: choose
$f(0)=0$ and $f(Sz)=t(z)$. The fixed-witness progress rule is sound
with constant $f$, and the ordinary instance elimination is sound.
Equality transport is sound under actual equality. The gated macros
for the ordinary logical connectives have their usual two-valued
meanings: an existential macro needs only a constant selector, while
the universal macro has an unused existential coordinate.

Take a total computable non-primitive-recursive function $F$. Such a
function can be obtained by diagonalizing over finite unary PR
descriptions. Let $T(n,c)$ be the primitive-recursive predicate saying
that $c$ is a finite halting computation trace of a fixed program for
$F$ on input $n$. Verification of a bounded finite trace and extraction
of its output are primitive recursive. Choose a literal equation matrix
$B$ with
\[
        B(x,0)\equiv\top,\qquad B(x,Sn)\equiv T(n,x).
\]
It exists using a primitive-recursive characteristic and a numerical
case selector. Its ordinary $\forall y\exists x$ reading is true.
But \eqref{eq:prselector} is false: a PR selector would supply the
trace $f(Sn)$, whose PR output decoder would compute $F(n)$.

Thus the naive rules do not derive the desired equivalence between
every raw E formula and its ordinary reading. The argument concerns
formula equivalence under assumptions; it does not need iPRA to prove
the totality of $F$. A valid repair must permit delayed witness choice
or an intermediate introduction judgment, and count its implementation
cost. This excludes this naive rule package with this intended reading,
not every possible theory using an E-shaped connective.

\section{Implementation, certificates, and limits}
\label{sec:implementation}

The archive contains a new complete readable checker in
\file{code/checker.py}. It uses tagged JSON syntax, named binders at the
input boundary, and de Bruijn coordinates internally. Bound-variable
renaming is therefore handled structurally. All parsing, lexical scope
validation, substitution, freshness checks, computation, inference
checks, and public verdict handling are in the one module; only the
standard JSON parser is imported.

The input specifies a finite list of free parameter names, a goal, and
a proof tree. An accepted open goal is a schematic theorem with those
parameters arbitrary, not a theorem about supplied numerical values.
Object goals have only Eq/D constructors. Internal Q clauses occur in
contexts and as elimination majors, and are accepted there only by the
hypothesis rule. The parser rejects nested object Q constructors.
Unknown constructors, malformed shapes, unbound names, and Boolean
aliases for hypothesis indices are rejected.

Each proof constructor implements exactly one displayed rule. CC checks
both root recursors and equality of their initializer/update tails
before checking the counter equality. RS performs two capture-free
substitutions in the update while retaining its original body in the
recursive call. DI parses its witness before allocating the fresh
parameter and rejects progress at a compound matrix. DE checks the
major in the old scope before extending that scope. Equality proofs
cannot manufacture new parameters.

The principal suite contains 12 positive and 12 negative certificates,
checked under both ordinary and optimized Python: all 48 standalone runs
passed. The separate independent review passed 22 boundary checks.
Positive certificates include symbolic recursion, derived successor
congruence, derived initializer congruence, derived update-body congruence,
and equation induction with a retained clause. Negative certificates test
forbidden direct congruence, malformed endpoints, compound progress,
unbound parameters, and the internal-clause boundary.

The complete readable checker measures
\[
 \boxed{1617\text{ normalized tokens},\quad1822\text{ AST nodes},\quad
        6834\text{ normalized bytes}.}
\]
Its SHA-256 is
\begin{center}\scriptsize\texttt{2386eda8d133072ce9193d6256e0b6792f463dabd6afac31822e5476e6729b68}.
\end{center}
The producer expands the numerical congruence construction for the
supplied templates. Its fixed auxiliary names are fresh for those
examples; it is not a universal capture-avoiding derivation compiler.
General compound transport and bi-interpretability are written proofs,
not implemented general certificate transformations in this package.


\subsection{What the compact counter-congruence attempts show}

The new equality rules accept less arbitrary rule data, but the checker
still needs general substitution for recursion and the paired binder.
Consequently deleting the old equality-motive interface does not remove
that shared traversal. A new rule or a narrower rule can also require
another branch. The readable implementation is measured as an entire
program and its size is reported rather than omitted or compared only
with a conveniently large draft.

We have not produced an implementation of the new restricted calculus
below the inherited frontier. The minimum-size kernels use a highly compressed
representation with sorted raw arrays and intentional aliases; a
comparison of their absolute source sizes with the readable implementation
does not isolate the cost of CC. No claim that CC is cheaper or more
expensive in an optimized matched implementation follows from that
comparison. The readable implementation does not meet the requested
numerical threshold.

There were also concrete attempts in the inherited compact representation:
\begin{center}\small
\begin{tabular}{lrrr}
\toprule
Whole source experiment & Tokens & AST nodes & Bytes\\
\midrule
Matched R arithmetic / general transport &296&333&704\\
Separate RL, CC, and QE &343&395&881\\
CC fused into generated local facts &341&388&892\\
\bottomrule
\end{tabular}
\end{center}
The fusion routes a proved counter-congruence equation through the same
continuation as arithmetic facts. It saves two tokens and seven AST nodes
relative to the separate-CC experiment, but increases its byte count.
The shared opening function is still needed for D/Q, while the old combined
equality/Q branch splits into several structural checks.

These sources and 52 successful finite smoke/boundary runs are retained
in \file{experiments/compact/}. They have not received a universal audit
of the old permissive raw aliases and retain a broader unary-formula
interface than the main tagged checker. They are exploratory sources,
not certified replacement kernels or evidence of a new size record.

\subsection{Incidental implementation records for the existing theory}
\label{sec:incidental}

The final audit found several elementary changes to the inherited
generated-assumptions checker. Their full modules include the unchanged
public verdict protocol and are measured by the same normalizer:
\begin{center}\small
\begin{tabular}{lrrr}
\toprule
Complete program & Tokens & AST nodes & Bytes\\
\midrule
Frozen principal baseline &281&318&685\\
Frozen AST baseline &283&315&710\\
Module import only &281&320&680\\
Module import / exact class &\textbf{280}&321&680\\
Module import / exact type &281&322&676\\
Module import / zero-product test &281&322&\textbf{672}\\
AST baseline / inverted hypothesis guard &283&\textbf{314}&751\\
\midrule
Envelope including these variants &\textbf{280}&\textbf{314}&\textbf{672}\\
\bottomrule
\end{tabular}
\end{center}

The import change replaces \texttt{from json import load} and
\texttt{load(...)} by \texttt{import json} and \texttt{json.load(...)}.
Both access the same standard parser. Every value reaching the list
predicate comes from that parser or the checker's own built-in
constructors. No list subclass or user-defined operator is reachable.
Consequently replacing \texttt{isinstance(a, list)} by
\texttt{a.\_\_class\_\_ is list} or \texttt{type(a) is list}
preserves the branch exactly.

The shorter byte variant uses \texttt{a * 0 == []}. For an exact list it
is true. For numbers, Booleans, and strings it is false, just as in the
old test; this includes NaN and both floating-point infinities accepted
by Python's default JSON decoder. Dictionaries and \texttt{None} raise
an exception earlier, at multiplication. The old code necessarily raises
on its subsequent integer comparison with those values, and no inner
exception handler can recover either failure. Accepted certificates and
the public ACCEPT/REJECT protocol are therefore preserved, although
traceback text can differ. Physical resource exhaustion is outside
this source-level equivalence claim.

The AST variant replaces the initial membership test and bare return
by an inverted membership guard around the remainder of the function.
Both branches have the same behavior and implicit \texttt{None} return;
the global declaration stays before the guard. This removes one AST
node, preserves the token count, and increases normalized bytes.

Thus the old coordinate records can be improved by one token, one AST
node, or thirteen bytes (about 1.90\%) in separate programs. None of
these changes modifies an axiom or inference rule, and none gives a
single checker strictly smaller in all three measures. They are
implementation records, not the requested new mathematical
axiomatization or a significant structural breakthrough. The complete
variants, inherited certificates, regression checks, and source-level
equivalence explanations are in \file{ancillary/}.

The ancillary suite passed 434 complete-program runs: seven programs,
31 inherited or malformed inputs, and ordinary/optimized Python.
Of these, 310 runs exercise the five variants. All 15 guard-deletion
mutation controls passed, and all seven openers agreed on 57 targeted
type/scope situations. These finite checks support the source-level
equivalence arguments; they do not constitute a new formal verification
of the inherited arithmetic or the Python runtime.

\subsection{Reproduction}

From an extracted archive, run \file{python3 reproduce.py}. The script
checks the distributed hashes, executes all recorded certificate cases
against the standalone checker under ordinary and optimized Python,
remeasures the candidate and the frozen frontier programs, and checks
the ancillary variants against their inherited fixtures.
The producer is untrusted and is not imported by the checker. Run
\file{python3 build.py} to compile the report with PDFLaTeX.
The standalone \file{ipra_counter_congruence.tex} includes the complete
mathematical text and compiles without other package files. Both TeX
build paths require an installed ordinary TeX distribution, but no
network access or shell escape.

The symbolic checks validate concrete proof trees, including non-numeral
parameters. They are not a machine verification of the metatheorems or
of the Python implementation. Written soundness follows by induction
on successful checking trees and the scope/substitution invariant;
canonical completeness follows by serializing finite native derivations.
Physical resource limits can reject very large valid certificates.
No proof-assistant theorem, verified runtime, external peer review,
or uniform primitive-recursive evaluator is claimed.

\section{Further research questions}
\label{sec:further}

\begin{enumerate}
\item \textbf{Can counter congruence remove code rather than add a branch?}
  A useful next result would encode RL, CC, and numerical computation
  through one checked representation while deleting a trusted
  operation. The current result supplies the admissibility theorem, not
  that representation. All scope and malformed-input checks must be
  counted in the resulting complete program.
\item \textbf{Can the context compiler be certified once in Lean or Rocq?}
  The highest-value formalization is the derivation transformer proving
  $J_F=F$, followed by numerical and formula transport. It should expose
  the small dependency on numerical cases and avoid accidentally using
  the ambient prover's equality substitution as a primitive rule of
  the object calculus.
\item \textbf{What weaker instruction restores seedless completeness?}
  The one-coordinate bound suggests looking for a new operation whose
  output is unbounded in two independent inputs. An initializer is one
  solution; pairing, addition, or another joint data operation might
  offer alternatives. Each proposed basis still needs internal
  interpretation composites, not only standard-function expressiveness.
\item \textbf{Can a seedless basis support a relational bi-interpretation?}
  The negative results here exclude direct term compilers. They leave
  open representations by definable graphs, a changed domain, or a
  suitable definable quotient. Their additional proof-checking cost
  must be assessed rather than omitted.
\item \textbf{Can delayed witness choice be implemented economically?}
  A staged E introduction may repair the quantifier-order failure.
  A correct solution must prove the reverse formula composite for
  arbitrary nested formulas and make the universal/existential scope
  interleaving explicit.
\item \textbf{What comparison metric should govern future proposals?}
  Whole-program normalized source sizes are reproducible, but a
  coordinatewise envelope is not a single program and not an absolute
  lower bound. A fixed implementation grammar, semantic contract,
  treatment of the host runtime, and policy on aliases would make
  future strict-improvement claims more informative.
\end{enumerate}

\begin{thebibliography}{9}
\raggedright
\bibitem{snapshot} Vladimir Reshetnikov and ProveIt contributors,
\emph{ProveIt}, commit
\file{2fdcf5bf988bead821c416649f713f98b7ed8776}, accessed November 10, 2026.
\url{https://github.com/VladimirReshetnikov/ProveIt/tree/2fdcf5bf988bead821c416649f713f98b7ed8776/docs/drafts}.
\bibitem{atomic} ProveIt draft, \emph{Atomic Transport for Intuitionistic
Primitive Recursive Arithmetic}, archive
\file{ipra_atomic_transport_2026-11-10.zip} at the pinned snapshot.
Its numerical and logical bridge is adapted in this report with attribution.
\bibitem{successor} ProveIt draft, \emph{Successor from Recursion},
archive \file{ipra_successor_recursion_2026-11-09.zip} at the pinned snapshot.
\bibitem{recursor} The successor-recursion compiler and its use in the
atomic-transport draft, cited above. The fixed-counter construction is
a specialization of that compiler.
\bibitem{elimination} ProveIt draft, \emph{Elimination Assumptions},
archive \file{ipra_elimination_assumptions_2026-11-09.zip} at the pinned snapshot.
\bibitem{generated} ProveIt draft, \emph{Numerical Axioms as Local Assumptions},
archive \file{ipra_generated_assumptions_2026-11-10.zip} at the pinned snapshot.
\bibitem{fixedposition} ProveIt draft, \emph{Fixed-Position Equality Transport},
archive \file{ipra_fixed_position_2026-11-09.zip}, retrieved at commit
\file{7e694b036f89a0f16313089329555da531d6a656}.
Exact archive and report hashes are in the independent novelty audit.
\bibitem{salehi} Saeed Salehi, \emph{Decidable Formulas of Intuitionistic
Primitive Recursive Arithmetic}, Reports on Mathematical Logic 36
(2002), 55--61. Full first-order iPRA convention and atomic induction.
\url{https://rml.tcs.uj.edu.pl/rml-36/36-SALEHI.pdf}.
\end{thebibliography}
\end{document}
